\section{The estimating problem}\label{sec:introduction}

The zeros of an entire function can be represented either by poles of
its logarithmic derivative or by point masses in the Laplacian of its
logarithmic modulus. For the Newman family these descriptions lead to
a positive, exponentially localized Green potential. They also isolate
a distinction that governs the analysis: an exact representation of a
zero signal, an accurate approximation to that signal, and an upper
bound excluding a zero are three different statements.

The field studied here is
\[
 H_\lambda(z)=\int_0^\infty \ee^{\lambda u^2}\Phi(u)\cos(zu)\dd u,
 \qquad
 H_0(z)=\frac18\xi\left(\frac12-\frac{\ii z}{2}\right),
\]
with the literal complete theta source \(\Phi\) specified in
Section~\ref{sec:source}. We use \(0\le\lambda\le1/2\) for the Green
geometry and the inherited canonical zero strip \(|\Im z_j|\le1\).
The arithmetic estimates are restricted to \(\lambda=0\).
Real symmetry means that a nonreal zero has an upper-half-plane
partner. All zero sums count multiplicity.

The main observable \(J_\lambda(X)\), defined in
Section~\ref{sec:green}, has the exact form
\[
 J_\lambda(X)=\sum_{\Im z_j>0}
 \atanh\frac{\sin(\pi\Im z_j/8)}
               {\cosh(\pi(X-\Re z_j)/8)}.
\]
Thus a zero with \(\Im z_j\ge d>0\) and
\(|\Re z_j-X|\le r\) forces an explicit positive margin. The desired
estimate has the opposite direction: a separately obtained upper
value must lie strictly below that margin. The potential does not
detect the multiplicity of a zero already on the real boundary.

Our development has four stages. First we establish the source,
resolvent, complete Green representation and exact prime subtraction.
Second we prove a quantitative stability obstruction on a specified
enlarged source class. Third we analyze quadratic approximation to a
pole-cancelled zeta field and prove, for fixed \(0<\theta<1/2\),
the residual floor
\[
 \liminf_{T\to\infty}\inf_A
 \frac1{|I_T|}\int_{I_T}\epsilon_A(X)^2\dd X
 \ge\frac{1/(2\theta)-11/12}{\zeta(5/2)^4}>0.
\]
Here a single bounded-coefficient Dirichlet polynomial of length at
most \((T/2)^\theta\) is used throughout a padded interval \(I_T\).
Finally we pass directly to the boundary logarithm, deriving a
one-sided error modulus from a safe interior anchor and a disk
growth bound. This last theorem allows real multiple zeros and does
not replace the measured logarithm by a harmonic or quadratic
surrogate.

The complete source and Green background are developed in the parent
manuscript and its assessed Green note \cite{OGParent,GreenStripProof}.
The source-class obstruction and quadratic floor are the branch
results \cite{OGConditioning,OGMollifier}, read with the correction and
assessment \cite{OGCorrection,OGAssessment}. The local logarithmic
theorem is \cite{OGBoundaryLog}. We give their derivations here in a
single notation. Classical complex-analysis ingredients and the
external mean-value theorem are attributed where used. We make no
literature-priority claim.

The two obstruction theorems have different domains. The first varies
the source away from theta and concerns stability of prescribed
observations. The second keeps the exact canonical field and varies
a restricted preconditioner. Neither gives a lower bound for the
canonical Green potential. The new local logarithmic theorem is an
acquisition result: it controls the error of retaining a signed
quantity whose favorable value is still unknown. Section~\ref{sec:tasks}
states precisely what remains to complete that passage.

